\section{Introduction}
Studies of high-multiplicity proton--proton (\pp) and proton--lead (\pPb) collisions have revealed that small collision systems exhibit signatures previously considered unique features of heavy-ion collisions. Some of these signatures, such as the enhanced production of strange hadrons~\cite{Naturepaper}, and collective flow~\cite{Ref:Flow1}~\cite{Ref:Flow2}, can be explained by the formation of a strongly interacting medium. Strangeness enhancement was one of the first proposed quark--gluon plasma (QGP) signatures~\cite{Koch:1986ud}, as the QGP-transition temperature is typically expected of being in the order of the strange quark mass, allowing for thermal production of strange quarks in a QGP. Collectivity was historically also expected to require thermalization and provide information on the equation of state~\cite{Teaney:2000cw}, while more recently it has been understood how collectivity can build up from kinetic equilibrium alone~\cite{Kurkela:2018xxd}. However, the formation of a medium in these small systems challenges current theoretical frameworks, because their initial small volumes imply lifetimes so short that it is unclear to what degree the systems can equilibrate (see Ref.~\cite{Kurkela:2018xxd} and references therein).

The observation of collective flow, as well as strangeness enhancement in particular, implies that pp collisions at LHC and RHIC energies can no longer be described as semi-incoherent sums of parton--parton collisions, an idea that has been central to most general-purpose quantum chromodynamics (QCD)-inspired Monte-Carlo event generators, such as PYTHIA~\cite{Sjostrand:2007gs} and Herwig~\cite{Herwig}. “Jet Universality'' is another long-standing idea for the phenomenological understanding of QCD assuming that, while the partonic processes vary with system and beam energy, the produced color fields and their hadronization are universal. In the context of Lund strings, this implies that the string tension and string-fragmentation parameters are the same regardless of collision system. For a recent discussion, we refer to Ref.~\cite{Torbjorn}. The discovery of strangeness enhancement scaling with the multiplicity~\cite{alice_newstrangeness}\cite{ALICE:2020nkc} violates the assumptions of jet universality, and for this reason QCD-inspired generators have incorporated additional phenomenological final-state pre-hadronization mechanisms, such as string percolation~\cite{Ref:SPerc}, color ropes~\cite{Ropes}, baryon junctions~\cite{Junctions} and/or new types of baryon-favored color reconnection~\cite{HerwigStrange}. Furthermore, charm fragmentation fractions in pp collisions also differ significantly from the values measured in $\textrm{e}^+\textrm{e}^-$ collisions~\cite{ALICE:2021dhb}. 


In contrast, QGP-inspired models include a system evolution with volume and multiplicity and so the qualitative features of both collective flow and strangeness enhancement are expected. EPOS-LHC is an event generator where the initial interactions lead to a two-phase state (core--corona) consisting of a dense core of QGP, and a diluted corona ~\cite{EPOS-LHC}. Strangeness enhancement in EPOS-LHC is due to a change in the relative contribution of the corona (low strangeness production from pomerons) and core (high thermal strangeness production from QGP) with multiplicity.

As can be seen from the previous discussion, strange-particle production
appears to be a powerful probe for QGP-like effects in small systems, and the origin of the
strangeness enhancement constitutes an important open question. Therefore, results on strangeness production in small systems can facilitate progress in the understanding of both QCD dynamics and hadronization, in addition to putting constraints on phenomenological models.

ALICE has previously reported that strangeness enhancement as a function of the average multiplicity at midrapidity does not depend on the collision system nor center-of-mass energy per nucleon--nucleon collision, ranging from pp to p--Pb and \snn = 2.76 TeV to \s =13 TeV~\cite{13tevvsmultpapers}. This universality implies a strong correlation between the underlying physics processes that drive both the enhancement and the multiplicity. In this article, we have tested if the observed universality can be broken by contrasting high-multiplicity events dominated by one or more hard scatterings, with events dominated by multiple softer interactions.

At high transverse momentum (\pt), strangeness production in hadronic collisions tends to originate from “hard'', perturbative QCD processes. These hadrons are either produced directly through flavor creation ($\rm XX \rightarrow \rm s\bar{s}$) or flavor excitation ($\rm{s}X \rightarrow \rm{s}X$), or indirectly as a result of radiation and/or hadronization associated with hard processes, e.g., through gluon splitting following the partonic evolution ($\rm g \rightarrow s\bar{s}$). In contrast, the production of “soft'', low-\pt strange hadrons ($\pt \leq 2~\gev$) is dominated by non-perturbative QCD processes, where novel QCD dynamics could be found. The final-state azimuthal topology is expected to reflect which of these QCD processes are primarily driving particle production for a given event. Events dominated by one or more hard scatterings will presumably lead to pronounced back-to-back jet structures, while events that contain several softer scatterings will result in isotropic distributions. We note that previous results obtained by ALICE, both utilizing the same event shape estimator~\cite{ALICE:2019dfi} and a similar one, the transverse sphericity~\cite{ALICE:2012cor}, found evidence for this behavior in the way that \meanpt would depend on the event topology. 

To identify the final-state azimuthal topology, we will here use a modified variant of the transverse spherocity (\SO) estimator proposed in Ref.~\cite{ALICE:2019dfi}, for being more sensitive to the underlying processes. The lower bound of \SO aims to distinguish “Jet-like'' events, which in this study are events characterized by an azimuthal topology similar to a pair of back-to-back jets (implying a tight clustering of particles with a difference in azimuthal angle $\Delta\phi \approx 0$ or $\pi$~\cite{ALICE:2019dfi}), which on average produces a larger number of high-\pt hadrons compared with the \SO--integrated distribution, thereby “hardening'' the \pt-differential spectra. Conversely, the upper bound of \SO selects “Isotropic'' events, defined by an azimuthal topology which is close to symmetric, with an absence of preferred direction. The hypothesis is that \SO, by contrasting events with these different topologies, can be used to control the degree of QGP-like effects, like strangeness enhancement and radial flow, in high-multiplicity \pp collisions~\cite{Antonio2}~\cite{Antonio3}.

In this article, we present the first results on strangeness production in high-multiplicity pp collisions at \sppt{13} as a function of the transverse spherocity. The results obtained for the light-flavor hadrons are presented as the sum of particles and anti-particles, explicitly as \pip+ \pim, \kap+ \kam, \KSTAR + \KSTARb, p + \pbar, \lmb+ \almb, \X +  \Ix where the exceptions are \kzero and \PHI. Hereinafter, the sum of particle and anti-particles will be referred to as $\pi$, K, \kzero, \KSTAR, p, \PHI,  \lmb, and $\Xi$, unless otherwise explicitly mentioned. 

The article is organized as follows. The ALICE main detectors used in this analysis are detailed in Section~\ref{Sec:ALICE}. Section~\ref{Sec:Mult} describes the high-multiplicity definitions used throughout this article. The transverse spherocity observable, along with caveats, is defined in Section~\ref{Sec:SOPT}. The details concerning particle identification (PID), yield extraction, and correction procedures are discussed in Section~\ref{Sec:Details}. The details regarding experimental corrections and systematic uncertainties are discussed in Sections~\ref{Sec:Corr} and~\ref{Sec:Sys}, respectively. The results are reported in Section~\ref{Sec:Results}, and finally the summary and conclusions of our study are discussed in Section~\ref{Sec:Conc}.



